% ============================================================================
% frontMatter.tex — dedicaces, remerciements, listes, introduction generale
% (numerotation romaine ; l'intro generale termine les pages liminaires)
% ============================================================================

\chapter*{Dedicaces}
\addcontentsline{toc}{chapter}{Dedicaces}

\vspace{2cm}
\begin{flushright}
\textit{
A nos chers parents, pour leur soutien indefectible et leurs sacrifices\\
tout au long de notre parcours academique,\\[0.4cm]
A nos professeurs, pour la qualite de leur enseignement\\
et la rigueur qu'ils nous ont transmise,\\[0.4cm]
A toute personne ayant contribue, de pres ou de loin,\\
a la reussite de ce projet de fin d'annee.
}
\end{flushright}

\newpage

\chapter*{Remerciements}
\addcontentsline{toc}{chapter}{Remerciements}

Nous tenons a exprimer notre profonde gratitude a toutes les personnes qui ont contribue,
de pres ou de loin, a la reussite de ce projet de fin d'annee.

Nous remercions chaleureusement notre encadrant professionnel, \textbf{M. Haitam BOUTAFI},
Ingenieur au sein de \textbf{HPS Group Morocco}, pour son encadrement de qualite, sa
disponibilite constante et ses conseils techniques avises qui ont guide la realisation
de ce projet d'un bout a l'autre.

Nous remercions egalement notre encadrant pedagogique, \acompleter{[Prenom NOM \textperiodcentered\ Grade]},
pour son suivi rigoureux et ses recommandations methodologiques tout au long de ce projet.

Nos remerciements s'adressent aussi a l'ensemble du corps professoral et administratif
de l'EMSI, pour la qualite de la formation dispensee durant ces cinq annees, ainsi qu'a
\textbf{HPS Group Morocco} pour nous avoir accueillis et permis de travailler sur une
problematique technique riche et formatrice.

Enfin, nous remercions nos familles et nos proches pour leur soutien moral sans faille.

\newpage

\tableofcontents
\newpage

\listoffigures
\addcontentsline{toc}{chapter}{Liste des figures}
\newpage

\listoftables
\addcontentsline{toc}{chapter}{Liste des tableaux}
\newpage

\chapter*{Liste des abreviations}
\addcontentsline{toc}{chapter}{Liste des abreviations}

\begin{table}[H]
\centering
\renewcommand{\arraystretch}{1.3}
\begin{tabularx}{\textwidth}{@{} l X @{}}
\toprule
\textbf{Abreviation} & \textbf{Signification} \\
\midrule
PFA & Projet de Fin d'Annee \\
\rowcolor{gray!6} PoC & Proof of Concept \\
ISO & International Organization for Standardization \\
\rowcolor{gray!6} MTI & Message Type Indicator (ISO 8583) \\
PAN & Primary Account Number \\
\rowcolor{gray!6} CDC & Change Data Capture \\
JDBC & Java Database Connectivity \\
\rowcolor{gray!6} SMT & Single Message Transform (Kafka Connect) \\
PK & Primary Key (cle primaire) \\
\rowcolor{gray!6} CRD & Custom Resource Definition (Kubernetes) \\
HA & Haute Disponibilite (High Availability) \\
\rowcolor{gray!6} PVC & Persistent Volume Claim \\
KRaft & Kafka Raft (metadata sans ZooKeeper) \\
\rowcolor{gray!6} AES & Advanced Encryption Standard \\
HMAC & Hash-based Message Authentication Code \\
\rowcolor{gray!6} TLS & Transport Layer Security \\
RBAC & Role-Based Access Control \\
\bottomrule
\end{tabularx}
\caption{Liste des abreviations utilisees}
\label{tab:abreviations}
\end{table}

\newpage

% ---- A partir d'ici : numerotation arabe, page 1 --------------------------
\clearpage
\pagenumbering{arabic}
\setcounter{page}{1}

\chapter*{Introduction generale}
\addcontentsline{toc}{chapter}{Introduction generale}

Le secteur des paiements electroniques connait une croissance soutenue, portee par la
digitalisation des services bancaires et la multiplication des canaux de paiement
(carte, mobile, sans contact). Cette croissance s'accompagne d'une exposition accrue
aux tentatives de fraude, dont la detection efficace constitue un enjeu majeur pour les
etablissements financiers et les operateurs de switching.

\textbf{HPS Group Morocco}, acteur reconnu de l'edition de solutions monetiques avec sa
plateforme \textit{PowerCARD}, traite quotidiennement un volume important de transactions
de paiement. Le systeme de detection de fraude actuellement en production repose sur un
traitement par lots (\textit{batch}) execute la nuit, ce qui introduit une latence de
plusieurs heures entre l'occurrence d'une transaction et l'identification d'une eventuelle
anomalie. Dans le contexte de la fraude bancaire, ou la fenetre d'action utile se mesure
en secondes, cette latence rend toute intervention corrective largement inoperante.

C'est dans ce contexte que s'inscrit notre projet de fin d'annee, qui consiste a concevoir
et a realiser une \textbf{preuve de concept (PoC)} d'une plateforme de traitement de
donnees en temps reel, capable d'ingerer, de valider, d'enrichir et de scorer des
transactions bancaires au fil de l'eau, en s'appuyant exclusivement sur un ecosysteme
\textbf{100\% open source} (Apache Kafka, Apache Flink, MinIO, PostgreSQL) deploye sur
une infrastructure Kubernetes orchestree par Terraform.

Ce rapport est organise en quatre chapitres. Le \textbf{premier chapitre} presente le
contexte general du projet : l'organisme d'accueil, l'etude de l'existant et la
methodologie de gestion de projet retenue. Le \textbf{deuxieme chapitre} est consacre a
l'analyse des besoins fonctionnels et non fonctionnels, ainsi qu'a la modelisation des cas
d'utilisation. Le \textbf{troisieme chapitre} detaille la conception du systeme, tant sur
le plan dynamique et statique que sur le plan architectural. Enfin, le \textbf{quatrieme
chapitre} presente la realisation du projet : l'environnement technique mis en place et
les interfaces graphiques temoignant du fonctionnement effectif du pipeline.

\newpage
